% Part: incompleteness
% Chapter: theories-computability
% Section: computably-axiomatizable

\documentclass[../../../include/open-logic-section]{subfiles}

\begin{document}

\olfileid{inc}{tcp}{cax}

\olsection{\printtoken{S}{axiomatizable} సిద్ధాంతాలు}

ఒక సిద్ధాంతం~$\Th{T}$కు గణనీయ స్వీకృతాల సమితి~$A$ ఉంటే, దాన్ని
\emph{\tetoken{స్వీకృతీకరించదగినది}{!!{axiomatizable}}} అంటాం. ($A$ అనేది $\Th{T}$కు స్వీకృతాల సమితి
అనడం అంటే $T = \Setabs{!A}{A \Proves !A}$.) సహజ సంఖ్యల యొక్క ప్రతి
``సమంజసమైన'' స్వీకృతీకరణకూ ఈ ధర్మం ఉంటుంది. ముఖ్యంగా, పరిమిత స్వీకృతాల
సమితి గల ప్రతి సిద్ధాంతమూ \tetoken{స్వీకృతీకరించదగినది}{!!{axiomatizable}}.

\begin{lem}
$\Th{T}$ \tetoken{స్వీకృతీకరించదగినది}{!!{axiomatizable}} అనుకుందాం. అప్పుడు $\Th{T}$
\tetoken{గణనీయంగా లెక్కించదగినది}{!!{computably
enumerable}}.
\end{lem}

\begin{proof}
$A$ అనేది $\Th{T}$కు గణనీయ స్వీకృతాల సమితి అని అనుకుందాం. $!A \in T$
అవుతుందో తెలుసుకోవడానికి, ఆ స్వీకృతాల నుంచి $!A$కు \tetoken{ఒక వ్యుత్పత్తి}{!!a{derivation}} కోసం
వెతకండి.

కొద్దిగా భిన్నంగా చెప్పాలంటే, $A$లో $!B_1$, \dots, $!B_k$ అనే పరిమిత
స్వీకృతాల జాబితీ, మొదటిస్థాయి తర్కంలో $(!B_1 \land \dots \land !B_k) \lif
!A$కు \tetoken{ఒక వ్యుత్పత్తి}{!!a{derivation}} ఉన్నప్పుడు, అప్పుడు మాత్రమే $!A$, $\Th{T}$లో ఉంటుంది.
కాని ``\dots జరిగేలా ఏదో \dots ఉంది'' అనే రూపంలో నిర్వచనం గల ప్రతి సమితీ,
రెండవ ``\dots'' గణనీయమైనదైతే, \tetoken{గ.లె.}{!!{c.e.}} అని మనకు ఇప్పటికే తెలుసు.
\end{proof}

\end{document}
